% Complete finite-pattern proof of Kamp's first supporting lemma for III.2.

\section{Kamp's basic past and future lemma}
\label{sec:annex-p09-kamp-iii2-lemma1}

\begin{annexitem}{annex:p9.kamp.iii2.lemma1}{thm:p09-kamp-iii2-lemma1}{Temporal expressions for basic finite interval patterns}
Use Definitions~\ref{def:p09-kamp-basic-expressibility}--\ref{def:p09-kamp-basic-word}.
The following argument proves the finite-chain rearrangement used in Kamp's
Lemma~1, and supplies both its past and future constructions explicitly
\cite[ch.~III, pp.~63--66, 69--70]{Kamp1968}.

\begin{proof}
Fix an arbitrary nonempty strict linear order and an arbitrary interpretation
of every unary predicate, giving the structure \(\mathcal M\).  All
renamings of bound variables below preserve their scopes and avoid the
endpoint names.

\paragraph{Word length and free variables.}
At stage one, either binary pattern has the three predicate occurrences
\((A,B,C)\), and exactly \(x,y\) free.  At the next past stage,
the new endpoint letter \(A\) and interval letter \(B\) precede all
the constituent's occurrences.  At the next future stage, the new interval
letter \(B\) and endpoint letter \(A\) follow all the constituent's
occurrences.  Each constructor therefore adds two occurrences, giving
\(2n+1\) at stage \(n\).  In the past constructor the constituent has
only \(z,y\) free; binding \(z\) leaves \(x,y\).  In the future
constructor the constituent has only \(x,z\) free, and binding \(z\)
again leaves \(x,y\).  The universal interval variable is fresh and bound.
Consequently the guarded existential closures have exactly their stated
single free variable.  Repeated predicate letters do not affect any of
these claims: the word records occurrences, not a set of distinct letters.

\paragraph{Binary patterns are finite chains on ordered endpoints.}
We prove simultaneously for the two binary classes that, whenever \(a<b\),
\[
 \mathcal M\models\phi[x\mapsto a,y\mapsto b]
 \quad\Longleftrightarrow\quad
 \text{there is a realization of \(\sigma(\phi)\)
       beginning at \(a\) and ending at \(b\)}.
 \tag{1}
\]
The condition \(a<b\) is part of this assertion.  It is not being
inserted into the binary source grammar.

At stage one, the endpoint conjuncts assert \(a\in A^{\mathcal M}\)
and \(b\in C^{\mathcal M}\).  The universal conjunct asserts that every
point strictly between them belongs to \(B^{\mathcal M}\).  These are
exactly the realization conditions for the sequence \((a,b)\), proving
both directions of (1) for both base classes.

Suppose (1) holds at stage \(n\).  A past-stage formula at stage
\(n+1\) has a constituent \(\theta\in\operatorname{Bp}_n(z,y)\).
If that formula holds at \((a,b)\), its existential supplies a point
\(c\) with \(a<c<b\); its other new conditions say
\(a\in A^{\mathcal M}\) and \((a,c)\subseteq B^{\mathcal M}\).
By induction, the truth of \(\theta\) with \(z\mapsto c,y\mapsto b\)
supplies a realization
\[
 c=a_1<a_2<\cdots<a_{n+1}=b
\]
of its word.  Prepending \(a=a_0\) gives a realization of the whole
word \((A,B,\sigma(\theta))\).  Conversely, a realization
\(a=a_0<\cdots<a_{n+1}=b\) of that whole word supplies the witness
\(c=a_1\).  The first point and first interval meet the new conditions,
and the remaining points realize \(\sigma(\theta)\).  The induction
hypothesis gives the truth of \(\theta\) at those endpoints, so the
past-stage formula holds.

For a future-stage formula at stage \(n+1\), the constituent is
\(\theta\in\operatorname{Bf}_n(x,z)\).  If the formula holds at
\((a,b)\), choose its witness \(c\) with \(a<c<b\).  Induction
provides a realization
\[
 a=a_0<a_1<\cdots<a_n=c
\]
of \(\sigma(\theta)\).  The new conditions are
\((c,b)\subseteq B^{\mathcal M}\) and \(b\in A^{\mathcal M}\);
append \(b=a_{n+1}\) to obtain the realization of
\((\sigma(\theta),B,A)\).  Conversely, given that whole realization,
use its penultimate point \(c=a_n\) as the existential witness.  The
initial sequence realizes the constituent, while the last interval and
point meet the remaining conditions.  Induction gives the constituent's
truth and hence the future-stage formula.  This completes the simultaneous
induction and proves (1) in full.

Every word of length \(2n+1\) has a pattern in each binary class: use
the three-letter base, recursively prepend its first two letters for a
past pattern, and recursively append its last two letters for a future
pattern.  By (1), any such two patterns with the same word agree whenever
their endpoints are ordered.  Thus the past and future ways of placing
the inner quantifiers describe the same finite-chain condition under the
guard required by the one-variable closures.
They also agree on unordered endpoint assignments: at stage one the same
word gives the same formula up to its bound-variable name, and at every
later stage each pattern's strict interior witness forces \(x<y\).
If that order fails both later-stage patterns are false.  This proves the
source's same-word equivalence for all assignments, as well as the guarded
form needed here.

\paragraph{The two guarded existential closures.}
If \(\phi\in\operatorname{Bp}_n(x,y)\), (1) immediately yields
\[
 \mathcal M\models\exists x\,(x<y\land\phi)[y\mapsto b]
 \quad\Longleftrightarrow\quad
 \text{\(\sigma(\phi)\) has a realization ending at \(b\)}.
 \tag{2}
\]
For the forward direction take the chosen \(a<b\) and apply (1).
For the reverse direction, the first point \(a_0\) of a realization
ending at \(b=a_n\) satisfies \(a_0<b\), because \(n\geq1\),
and (1) verifies the matrix at those endpoints.

Similarly, if \(\psi\in\operatorname{Bf}_n(x,y)\),
\[
 \mathcal M\models\exists y\,(x<y\land\psi)[x\mapsto a]
 \quad\Longleftrightarrow\quad
 \text{\(\sigma(\psi)\) has a realization beginning at \(a\)}.
 \tag{3}
\]
A witness for the existential and (1) prove the forward direction.  In
the reverse direction the last point \(a_n\) of the realization is
strictly later than \(a=a_0\); choose it as the witness and apply (1).
Neither argument creates an endpoint or an interior point absent from the
order.

\paragraph{Explicit temporal construction for the past closure.}
Let \(\sigma=(A_0,B_1,A_1,\ldots,B_n,A_n)\), and define \(L_i\)
as in the lemma.  We prove by induction on \(i\) that
\(D(L_i)[w/t_0]\), evaluated at \(w=b\), holds exactly when the
initial word \((A_0,B_1,A_1,\ldots,B_i,A_i)\) has a realization
ending at \(b\).  For this induction alone, the initial word with
\(i=0\) is realized by the single point \(b\) when
\(b\in A_0^{\mathcal M}\).  This is a semantic induction seed,
not an extra basic-pattern class.

For \(i=0\), the definition of \(D\) gives
\(D(L_0)[w/t_0]=A_0(w)\), establishing the assertion.  For
\(1\leq i\leq n\), the conjunction and strict-since clauses of \(D\)
give, up to fresh bound-variable names,
\[
\begin{split}
 D(L_i)[w/t_0]\ \longleftrightarrow\ 
 &A_i(w)\land\exists z\Bigl(z<w\land D(L_{i-1})[z/t_0]\\
 &{}\land\forall u\bigl((z<u\land u<w)
                            \longrightarrow B_i(u)\bigr)\Bigr).
\end{split}
\tag{4}
\]
If (4) holds at \(b\), take its witness \(c<b\).  Induction gives
a realization of the shorter word ending at \(c\).  Append \(b\):
the new gap satisfies \(B_i\) and the new endpoint satisfies \(A_i\).
Conversely, take a realization of the longer word ending at \(b\).
Its penultimate point \(c\) realizes the shorter word with its preceding
points.  Induction gives
\(\mathcal M\models D(L_{i-1})[t_0\mapsto c]\), and the last gap and
endpoint supply the other conjuncts of (4).  This proves the assertion
for \(i\), and hence for \(i=n\).  Combining it with (2) gives the
past equivalence claimed in the lemma.

\paragraph{Explicit temporal construction for the future closure.}
Define \(R_i\) as in the lemma.  By decreasing induction on \(i\),
\(D(R_i)[w/t_0]\) holds at \(w=a\) exactly when the terminal word
\((A_i,B_{i+1},A_{i+1},\ldots,B_n,A_n)\) has a realization
beginning at \(a\).  For \(i=n\), the terminal word has one letter,
and \(D(R_n)[w/t_0]=A_n(w)\) proves the assertion.

For \(0\leq i<n\), the conjunction and strict-until clauses give
\[
\begin{split}
 D(R_i)[w/t_0]\ \longleftrightarrow\ 
 &A_i(w)\land\exists z\Bigl(w<z\land D(R_{i+1})[z/t_0]\\
 &{}\land\forall u\bigl((w<u\land u<z)
                         \longrightarrow B_{i+1}(u)\bigr)\Bigr).
\end{split}
\tag{5}
\]
If this holds at \(a\), the witness \(c>a\) begins a realization
of the shorter terminal word by induction.  Prepending \(a\) gives the
whole realization: the first point meets \(A_i\) and the first interval
meets \(B_{i+1}\).  Conversely, the second point \(c\) of a whole
realization begins the required shorter realization.  Induction verifies
\(\mathcal M\models D(R_{i+1})[t_0\mapsto c]\), and the first interval
and endpoint verify (5).
The assertion follows down to \(i=0\).  Combining it with (3) proves
the future equivalence.

\paragraph{Endpoint cases, effectiveness and exact scope.}
All existential inequalities in (1)--(5) are strict.  A missing first or
last witness makes the corresponding closure false.  Universal interval
conditions are evaluated on open intervals and are true when those
intervals are empty.  Thus adjacent points, endpoints, finite orders and
orders with gaps all obey the same proof; no density or completeness
premise entered any step.

Read the finite occurrence word from the input pattern and perform exactly
the indicated finite recurrence for \(L_n\) or \(R_0\).  Each recurrence
uses the preceding temporal formula once and introduces only conjunction,
the appropriate strict temporal operator, and proposition letters already
corresponding to input predicates.  The definition of \(D\) then gives
an \(L\)-formula with the sole free current-point variable; the final
capture-avoiding renaming yields exactly the original free endpoint.

Kamp's proof reorders the existential witnesses into nested two-endpoint
links and then uses predicate substitution to express those links.  The
finite-chain equivalence (1) and the two recurrences (4)--(5) provide that
rearrangement and composition explicitly, without relying on an unproved
reordering assertion or leaving the future case implicit.  This proves
Lemma~1 over \(\mathsf{Lin}\); no conclusion about all first-order
regional formulas or the remaining proof of III.2 follows from this lemma
alone.
\end{proof}
\end{annexitem}
